\clearpage
\begin{proofbox}
\textbf{方法一：两个横坐标增量，用平方差建立关系}

沿用结论中的 $C,P,Q,M,N$。记 $D=1/a^2-m^2/b^2\ne0$、$m\ne0$，目标是证明 $M,N$ 在同一条 $ux/a^2+vy/b^2=t$ 上，从而平行于 $\ell_Q:ux/a^2+vy/b^2=1$。
\begin{center}
\begin{tikzpicture}[x=1cm,y=1cm,font=\small]
\draw[gray!45,rounded corners=3pt,fill=white] (0,0) rectangle (13.8,9.2);
\node[anchor=west,font=\small\bfseries] at (.4,8.7) {方法一：从同一个 P 出发，比较两个带符号的横增量};
\begin{scope}[shift={(1.25,4.15)},scale=1.25]\CNBase

\draw[cnhelp] (P)--(2,1.6) (M)--({10/7},1.6) (N)--({26/7},1.6);
\draw[-Stealth] (2,1.6)--({10/7},1.6);
\draw[-Stealth] (2,1.6)--({26/7},1.6);
\node[above] at (1.5,1.9) {$r$}; \node[above] at (3,1.9) {$s$};
\end{scope}
\node[anchor=west] at (8.1,7) {$r=x_M-u$};
\node[anchor=west] at (8.1,5.85) {$s=x_N-u$};
\node[anchor=west] at (8.1,4.6) {$x_M-x_N=r-s$};
\node[anchor=west] at (8.1,3.3) {$y_M-y_N=m(r+s)$};
\node[anchor=west] at (8.1,1.9) {图取 $r=-4/7$};
\node[anchor=west] at (8.1,.9) {$s=12/7,\ m=2$};

\end{tikzpicture}

\end{center}
\textbf{第1步：} 只设两个横增量：
\[
M=(u+r,v+mr),\qquad N=(u+s,v-ms),\quad r,s\ne0.
\]
\textbf{第2、3步：} 分别把 $M,N$ 代入双曲线，减去 $u^2/a^2-v^2/b^2=1$，展开平方并提取增量，得到
\[
\begin{aligned}
r\left(2u/a^2-2vm/b^2+Dr\right)&=0,\\
s\left(2u/a^2+2vm/b^2+Ds\right)&=0.
\end{aligned}
\]
因为 $r,s\ne0$（第二交点不是旧点 $P$），约去它们后分别有
\begin{align}
Dr&=-\frac{2u}{a^2}+\frac{2vm}{b^2},\tag{1}\\
Ds&=-\frac{2u}{a^2}-\frac{2vm}{b^2}.\tag{2}
\end{align}
\end{proofbox}
\clearpage
\begin{proofbox}
\textbf{方法一（续）：把两种增量组合成弦的方向}
\begin{description}[style=nextline,leftmargin=2.8em,labelsep=0.5em,itemsep=0.48em]
  \item[\textbf{第 4 步：先减后加，正好把两种量分开。}]
  用 (1) 减 (2)，再用 (1) 加 (2)：
  \[
   \boxed{D(r-s)=\frac{4vm}{b^2}},\qquad
   \boxed{D(r+s)=-\frac{4u}{a^2}}.\tag{3}
  \]
  为什么要加减？因为 $M,N$ 的横坐标差是 $r-s$，纵坐标差是 $m(r+s)$。

  \item[\textbf{第 5 步：不求斜率，也能判断平行。}]
  由坐标表示，
  \[
   x_M-x_N=r-s,\qquad y_M-y_N=m(r+s).
  \]
  把 (3) 代入下式的两项，发现恰好抵消：
  \[
  \begin{aligned}
  D\left[\frac{u(x_M-x_N)}{a^2}
    +\frac{v(y_M-y_N)}{b^2}\right]
    &=\frac{u}{a^2}\frac{4vm}{b^2}
      +\frac{vm}{b^2}\left(-\frac{4u}{a^2}\right)\\
    &=0.
  \end{aligned}
  \]
  $D\ne0$，于是
  \[
    \frac{ux_M}{a^2}+\frac{vy_M}{b^2}
    =\frac{ux_N}{a^2}+\frac{vy_N}{b^2}.
  \]
  所以 $M,N$ 在同一条 $\frac{ux}{a^2}+\frac{vy}{b^2}=t$ 上，与 $Q$ 点切线 $\ell_Q$ 平行。
\end{description}

\end{proofbox}
\clearpage
\begin{proofbox}
\textbf{\textcolor{CProof}{方法二：最快（一次求出“第二个交点”的增量）}}

\begin{center}
\begin{tikzpicture}[x=1cm,y=1cm,font=\small]
\draw[gray!45,rounded corners=3pt,fill=white] (0,0) rectangle (13.8,8.7);
\node[anchor=west,font=\small\bfseries] at (.4,8.2) {方法二：两交点相减，剩下固定的切线方向};
\begin{scope}[shift={(1.25,4.15)},scale=1.25]\CNBase

\draw[cnchord,-Stealth,line width=2pt] (M)--(N);
\draw[cntangent,-Stealth,line width=1.5pt] (1.5,-.5)--(2.7,-1.7);
\end{scope}
\node[anchor=west] at (8.1,7) {$R_k=P+h(1,k)$};
\node[anchor=west] at (8.1,5.7) {取 $k=m,-m$ 后相减};
\node[anchor=west] at (8.1,4.35) {$M-N=\dfrac{4m}{D}\,\boldsymbol d$};
\node[anchor=west] at (8.1,2.95) {$\boldsymbol d=\left(\dfrac v{b^2},-\dfrac u{a^2}\right)$};
\node[anchor=west] at (8.1,1.35) {方向成比例即可平行};

\end{tikzpicture}

\end{center}

\textbf{想法：} 先记住这一条简单的结果。直线 $y-v=k(x-u)$ 与双曲线交于 $P$ 和另一个点 $R_k$，令 $R_k=(u+h,v+kh)$，代入并减去 $P$ 的方程：
\[
 h\left[\left(\frac1{a^2}-\frac{k^2}{b^2}\right)h
       +2\left(\frac{u}{a^2}-\frac{vk}{b^2}\right)\right]=0.
\]
排除 $h=0$（对应旧点 $P$），得到
\[
 \boxed{R_k=P-\frac{2\left(u/a^2-vk/b^2\right)}{1/a^2-k^2/b^2}(1,k)}.
\]
把 $k=m$ 和 $k=-m$ 分别代入，即得
\[
\begin{aligned}
 M-N
 &=\left(\frac{4vm}{b^2D},-\frac{4um}{a^2D}\right)\\
 &=\frac{4m}{D}\left(\frac{v}{b^2},-\frac{u}{a^2}\right).
\end{aligned}
\]
而 $\left(\frac{v}{b^2},-\frac{u}{a^2}\right)$ 的方向沿着 $\ell_Q$，因为顺着这个方向走时，$\frac{ux}{a^2}+\frac{vy}{b^2}$ 不变。所以 $MN\parallel\ell_Q$。

\textbf{考场最省力之处：} 只对称替换 $m\mapsto-m$，直接把割线斜率 $m$ 从最终方向中消掉；不必把 $M,N$ 两套坐标全部写出来。

\end{proofbox}
\clearpage
\begin{proofbox}
\textbf{\textcolor{CProof}{方法三：常规法（分别联立直线，用韦达定理）}}

\begin{center}
\begin{tikzpicture}[x=1cm,y=1cm,font=\small]
\draw[gray!45,rounded corners=3pt,fill=white] (0,0) rectangle (13.8,8.7);
\node[anchor=west,font=\small\bfseries] at (.4,8.2) {方法三：同一个已知交点 P，对应两组中的同一个根 u};
\begin{scope}[shift={(1.25,4.15)},scale=1.25]\CNBase
\end{scope}
\node[anchor=west] at (8.1,7) {$PM$ 的横坐标根：};
\draw[-Stealth,gray] (8,5.7)--(13.2,5.7);
\fill ({8+1.25*10/7},5.7) circle(1.8pt) node[above left=5pt] {$x_M$};
\fill (10.5,5.7) circle(1.8pt) node[below right=5pt] {$u$};
\node[anchor=west] at (8.1,4.3) {$PN$ 的横坐标根：};
\draw[-Stealth,gray] (8,3)--(13.2,3);
\fill (10.5,3) circle(1.8pt) node[below=5pt] {$u$};
\fill ({8+1.25*26/7},3) circle(1.8pt) node[above=5pt] {$x_N$};
\node[anchor=west] at (8.1,1.4) {每组都含已知根 $u$};
\node[anchor=west] at (8.1,.65) {右侧仅表示根的对应};

\end{tikzpicture}

\end{center}

\begin{description}[style=nextline,leftmargin=2.8em,labelsep=0.5em,itemsep=0.45em]
 \item[\textbf{第 1 步：求第一条割线的两交点横坐标之和。}]
 直线 $PM$ 是 $y=mx+(v-mu)$。代入双曲线：
 \[
 Dx^2-\frac{2m(v-mu)}{b^2}x
       -\frac{(v-mu)^2}{b^2}-1=0.
 \]
 已知一个根是 $u$，另一个根是 $x_M$。由韦达定理：
 \[
  \boxed{x_M+u=\frac{2m(v-mu)}{b^2D}}.\tag{4}
 \]

 \item[\textbf{第 2 步：第二条割线同理。}]
 直线 $PN$ 是 $y=-mx+(v+mu)$，与双曲线联立，其两根是 $u,x_N$：
 \[
 \boxed{x_N+u=-\frac{2m(v+mu)}{b^2D}}.\tag{5}
 \]

 \item[\textbf{第 3 步：横坐标相减、相加。}]
 由 (4)、(5)，
 \[
 x_M-x_N=\frac{4vm}{b^2D},\qquad
 x_M+x_N=-2u-\frac{4m^2u}{b^2D}.
 \]

 \item[\textbf{第 4 步：求纵坐标差，消去 $m^2$。}]
 因为 $y_M=v+m(x_M-u)$、$y_N=v-m(x_N-u)$，
 \[
 \begin{aligned}
 y_M-y_N
 &=m(x_M+x_N-2u)\\
 &=m\left(-4u-\frac{4m^2u}{b^2D}\right)\\
 &=-4um\frac{D+m^2/b^2}{D}
 =-\frac{4um}{a^2D}.
 \end{aligned}
 \]
 最后一步使用了 $D+m^2/b^2=1/a^2$。

 \item[\textbf{第 5 步：写出弦的方向。}]
 因为 $M\ne N$ 且 $m,D\ne0$，上面两个坐标差不能同时为零。它们与 $\left(v/b^2,-u/a^2\right)$ 成相同比例，因此弦 $MN$ 与 $\ell_Q$ 平行。
 \end{description}

\medskip
\textbf{\textcolor{CProof}{补充：$v=0$ 为什么不能直接用斜率？}}

$v=0$ 时，(3) 给出 $r-s=0$，即 $x_M=x_N$，所以 $MN$ 竖直；此时 $P=(\pm a,0)$，$Q=P$ 点切线是 $x=\pm a$，也竖直。\textbf{平行结论照样成立；只是斜率不存在。}

\textbf{\textcolor{CProof}{三个证明共同说明的事实：}} $m$ 只能改变 $M,N$ 在哪儿，改变不了 $MN$ 的方向。证明并没有假定 $M$ 与 $N$ 是对称点，也没有提前猜出弦所在的具体直线。
\end{proofbox}
